\section{Segurança no Ambiente de Cloud e Contêineres}


O uso de computação em nuvem e a criação de cada vez mais aplicações \emph{cloud-native}, apoiado pela crescente preferência por modelos de arquitetura distribuída, principalmente os microsserviços são suportadas por aplicações por meio da transição de máquinas virtuais baseadas em hipervisores\footnote{Um hipervisor é um programa que acrescenta recursos de computação, como processamento, memória e armazenamento, e os redireciona entre máquinas virtuais (VMs)} para a virtualização em nível de sistema operacional via contêineres.

Enquanto na virtualização tradicional cada instância executa um sistema operacional convidado completo sobre um hipervisor sobre o hardware físico, os contêineres compartilham diretamente o mesmo \emph{kernel} do sistema hospedeiro \cite{sultan2019container}. Embora essa arquitetura reduza drasticamente o \emph{overhead} de memória e ciclos de processamento, ela aumenta a possibilidade de ataques em um sistema que contenha ambientes de múltiplos inquilinos (nesse caso, os contêineres). Uma vulnerabilidade na interface de chamadas de sistema ou uma falha lógica no \emph{kernel} compartilhado pode ser explorada para romper as da instância, resultando em ataques de escape de contêiner, nos quais um processo malicioso obtém acesso irrestrito ao sistema hospedeiro \cite{souppaya2017nist}.

Para controlar os privilégios desses contêineres e hipervisores, é necessário realizar um isolamento em nível de sistema operacional, isso é feito por meio de  primitivas de contenção do \emph{kernel}. Por meio de \emph{namespaces}, é possível isolar a visibilidade particionando recursos globais como árvores de identificadores de processos (\emph{PID}), pilhas de rede, pontos de montagem e identificadores de usuários, de modo que cada contêiner enxergue apenas seu próprio escopo operacional \cite{rice2020container}. Também é possível limitar o uso de recursos físicos de uma máquina por meio dos grupos de controle. Eles atuam na contabilização e limitação do consumo de recursos físicos como tempo de CPU, memória RAM e taxa de entrada/saída, prevenindo que uma instância comprometida execute ataques de negação de serviço por exaustão de hardware no nó físico.

Contudo, como os \emph{namespaces} e groupos de controle não filtram a totalidade das interações com o kernel, a defesa em profundidade exige a adoção de contêineres sem privilégios de superusuário (\emph{rootless containers}), a remoção de privilégios administrativos mediante a fragmentação de \emph{Linux Capabilities} (descartando permissões críticas como $cap\_sys\_admin$ \footnote{funciona essencialmente como um equivalente ao acesso tradicional de superusuário (\emph{root}) para uma vasta gama de operações do sistema.}), a restrição de chamadas de sistema via perfis Seccomp (\emph{Secure Computing Mode}) e o confinamento por Controle de Acesso Obrigatório \cite{rice2020container}.

A segurança em ambientes conteinerizados também é fortemente impactada por vulnerabilidades nas imagens desses e na cadeia de suprimentos de software. Imagens de contêineres são frequentemente derivadas de repositórios públicos não auditados, muitas dessas imagens contêm bibliotecas de software desatualizadas, pacotes com falhas de segurança conhecidas não corrigidas ou até software malicioso \cite{combe2016docker}. A mitigação desse risco requer varreduras automatizadas de vulnerabilidades diretamente nas esteiras de integração contínua (CI/CD), a assinatura criptográfica de artefatos para validação de proveniência no \emph{registry} e a adoção de imagens base mínimas ou \emph{distroless}. Essas imagens teiram apenas a aplicação contida no contêiner e suas dependências, eliminando interpretadores \emph{shell}, e gerenciadores de pacotes, que podem ser instrumentalizados por atacantes \cite{souppaya2017nist}.

Como abordado anteriormente, atualmente a maior parte dos sistemas de grande e médio porte opta por uma arquitetura de microsserviços, que expande exponencialmente o volume de tráfego interno entre contêineres, denominado tráfego Leste-Oeste . Por padrão, orquestradores priorizam
a conectividade plana (\emph{flat networking}), na qual qualquer contêiner (ou \emph{Pod}, conjunto de um ou mais contêineres) pode se comunicar livremente com outro endereço IP virtual do \emph{cluster} sem restrições nativas de camada de rede \cite{martin2021hacking}. Essa topologia facilita a movimentação de um invasor após o comprometimento inicial de um único microsserviço exposto. A neutralização desse vetor fundamenta-se na arquitetura \emph{Zero Trust}, implementada por meio de microssegmentação lógica via políticas de rede do tipo negação por padrão (\emph{default-deny}), criptografia e autenticação mútua na camada de transporte  gerenciada por malhas de serviço, e inspeção de pacotes de alta performance diretamente no espaço de kernel utilizando Extended Berkeley Packet Filter (eBPF) \cite{martin2021hacking}.

A criticidade prática desses vetores de ataque e da contenção de recursos evidenciou-se em incidentes recentes na infraestrutura do \emph{GitHub Actions}, cujos ambientes de integração contínua operam sobre \emph{runners} efêmeros e contêineres isolados em nuvem. No âmbito da cadeia de suprimentos e isolamento de privilégios, o comprometimento da ação \emph{tj-actions/changed-files} (CVE-2025-30066), em março de 2025, demonstrou como a ausência de restrições estritas de chamadas de sistema permite que códigos maliciosos inspecionem diretamente a memória do processo hospedeiro (\emph{/proc/PID/mem}) para extrair credenciais de \emph{registries} e chaves de nuvem em mais de 23 mil repositórios \cite{cisa2025tjactions}. De forma análoga, a campanha \emph{GhostAction}, documentada entre 2025 e 2026, explorou a ausência de microssegmentação e filtragem de tráfego de saída em contêineres de \emph{build} para obter mais de 3.300 segredos corporativos em cerca de 30 mil repositórios \cite{gitguardian2026ghostaction}. Adicionalmente, do ponto de vista da disponibilidade e do isolamento de recursos via \emph{cgroups}, a ocorrência de 48 interrupções críticas na plataforma ao longo de doze meses ilustra o impacto de falhas em cascata quando o particionamento de capacidade em ambientes de múltiplos inquilinos acaba falhando \cite{githubstatus2026}.